\documentclass[11pt]{article}

\usepackage[margin=1in]{geometry}
\usepackage{amsmath}
\usepackage{amssymb}
\usepackage{booktabs}
\usepackage{array}
\usepackage{xcolor}
\usepackage[hidelinks]{hyperref}

\title{Projected Functional Dependencies and Dependency Preservation}
\author{}
\date{}

\newcommand{\fd}{\rightarrow}
\newcommand{\code}[1]{\texttt{#1}}
\newcommand{\projlabel}[1]{\text{\small\bfseries #1}}
\setlength{\jot}{4pt}

\begin{document}

\maketitle

\setcounter{section}{4}

\section{Projecting the Functional Dependencies}

The minimal basis from that document is used as $F_{min}$. The two references to
\code{Region} are temporarily distinguished as \code{request\_region\_id} and
\code{server\_region\_id}, both are stored as \code{region\_id} in their respective tables as mentioned in the FD and Minimal Basis Document

\subsection{Projection procedure}

For each relation we:

\begin{enumerate}
    \item Choose a non-empty attribute set $C \subseteq S$.
    \item Calculate $\{C\}^{+}$ using $F_{min}$.
    \item Restrict the closure to $S$ using $\{C\}^{+} \cap S$.
    \item Record the non-trivial FDs and simplify them to a minimal basis.
\end{enumerate}

In principle every non-empty subset is considered. The first few projections are shown in
more detail. The remaining relations use the same procedure:

\[
\boxed{\text{Relation} \;\longrightarrow\; \text{Closure} \;\longrightarrow\;
\text{Intersection} \;\longrightarrow\; \text{Projected basis}}
\]

\subsection{Worked projection: Region}

\begin{align*}
\projlabel{Relation:} &\quad R = \{region\_id, region\_name, country\},\\
\projlabel{Closure:} &\quad \{region\_id\}^{+}
= \{region\_id, region\_name, country\},\\
\projlabel{Intersection:} &\quad \{region\_id\}^{+} \cap R = R.
\end{align*}

\begin{center}
    \textbf{Projected minimal basis:}
\end{center}
\[
region\_id \fd region\_name,
\qquad
region\_id \fd country.
\]

\subsection{Worked projection: GPUServer}

\begin{align*}
\projlabel{Relation:} &\quad GS = \{machine\_id, cpu\_capacity, memory\_capacity, gpu\_capacity,\\
&\qquad\qquad server\_region\_id, gpu\_type\_name\}.
\end{align*}

Using $F_{min}$:
\begin{align*}
\projlabel{Closure:} &\quad \{machine\_id\}^{+}
= \{machine\_id, cpu\_capacity, memory\_capacity, gpu\_capacity,\\
&\qquad\qquad server\_region\_id, gpu\_type\_name, memory\_gb\},\\
\projlabel{Intersection:} &\quad \{machine\_id\}^{+} \cap GS = GS.
\end{align*}

\begin{center}
    \textbf{Projected minimal basis:}
\end{center}
\begin{align*}
machine\_id &\fd cpu\_capacity &
 machine\_id &\fd memory\_capacity\\
machine\_id &\fd gpu\_capacity &
 machine\_id &\fd server\_region\_id\\
machine\_id &\fd gpu\_type\_name.&&
\end{align*}

Note that \code{server\_region\_id} is stored as \code{region\_id} in the table.

\subsection{Worked projection: MachineMetric}

\begin{align*}
\projlabel{Relation:} &\quad MM = \{machine\_id, metric\_id, start\_time, end\_time,\\
&\qquad\qquad worker\_load, gpu\_utilization, cpu\_utilization\}.
\end{align*}

First we check each part of the weak entity identifier:
\begin{align*}
\projlabel{Single-attribute closures:} &\quad
\{machine\_id\}^{+} \cap MM = \{machine\_id\},\\
&\quad \{metric\_id\}^{+} \cap MM = \{metric\_id\}.
\end{align*}

Neither alone determines the metric row. Thus:
\begin{align*}
\projlabel{Closure:} &\quad \{machine\_id, metric\_id\}^{+}\\
&\quad = \{machine\_id, metric\_id, cpu\_capacity, memory\_capacity,\\
&\qquad gpu\_capacity, server\_region\_id, gpu\_type\_name, memory\_gb,\\
&\qquad start\_time, end\_time, worker\_load, gpu\_utilization, cpu\_utilization\},\\
\projlabel{Intersection:} &\quad \{machine\_id, metric\_id\}^{+} \cap MM = MM.
\end{align*}

\begin{center}
    \textbf{Projected minimal basis:}
\end{center}
\begin{align*}
\{machine\_id, metric\_id\} &\fd start\_time &
\{machine\_id, metric\_id\} &\fd end\_time\\
\{machine\_id, metric\_id\} &\fd worker\_load &
\{machine\_id, metric\_id\} &\fd gpu\_utilization\\
\{machine\_id, metric\_id\} &\fd cpu\_utilization.&&
\end{align*}

This confirms that both attributes are required in the weak entity's full identifier.

\subsection{Worked projection: Hosts}

\begin{align*}
\projlabel{Relation:} &\quad H = \{machine\_id, model\_id\},\\
\projlabel{Closures:} &\quad \{machine\_id\}^{+} \cap H = \{machine\_id\},\\
&\quad \{model\_id\}^{+} \cap H = \{model\_id\},\\
\projlabel{Intersection:} &\quad \{machine\_id, model\_id\}^{+} \cap H = H.
\end{align*}

\noindent\textbf{Projected minimal basis:} No non-trivial FDs. The final intersection
produces only trivial dependencies.

\subsection{Remaining relation projections}

The remaining relations follow the same procedure as shown above.

\paragraph{GPUType.}
\begin{align*}
\projlabel{Relation:} &\quad G = \{gpu\_type\_name\},\\
\projlabel{Closure:} &\quad \{gpu\_type\_name\}^{+}
= \{gpu\_type\_name, memory\_gb\},\\
\projlabel{Intersection:} &\quad \{gpu\_type\_name\}^{+} \cap G = G.
\end{align*}
\noindent\textbf{Projected minimal basis:} No non-trivial FDs.

\paragraph{SpecificGPUType.}
\begin{align*}
\projlabel{Relation:} &\quad SG = \{gpu\_type\_name, memory\_gb\},\\
\projlabel{Closure:} &\quad \{gpu\_type\_name\}^{+}
= \{gpu\_type\_name, memory\_gb\},\\
\projlabel{Intersection:} &\quad \{gpu\_type\_name\}^{+} \cap SG = SG.
\end{align*}
\begin{center}
    \textbf{Projected minimal basis:}
\end{center}
\[
gpu\_type\_name \fd memory\_gb.
\]

\paragraph{AIModel.}
\begin{align*}
\projlabel{Relation:} &\quad M = \{model\_id, model\_name\},\\
\projlabel{Closure:} &\quad \{model\_id\}^{+}
= \{model\_id, model\_name\},\\
\projlabel{Intersection:} &\quad \{model\_id\}^{+} \cap M = M.
\end{align*}

\begin{center}
    \textbf{Projected minimal basis:}
\end{center}
\[
model\_id \fd model\_name.
\]

\paragraph{InferenceRequest.}
\begin{align*}
\projlabel{Relation:} &\quad Q = \{request\_id, submitted\_at, input\_size, status,\\
&\qquad\qquad response\_time\_ms, request\_region\_id, model\_id, machine\_id\}.
\end{align*}

The closure of \code{request\_id} also reaches attributes outside $Q$ through
\code{model\_id} and \code{machine\_id}; only attributes in $Q$ are kept during projection.
\begin{align*}
\projlabel{Closure:} &\quad \{request\_id\}^{+}\\
&\quad = \{request\_id, submitted\_at, input\_size, status, response\_time\_ms,\\
&\qquad request\_region\_id, model\_id, machine\_id, model\_name,\\
&\qquad cpu\_capacity, memory\_capacity, gpu\_capacity, server\_region\_id,\\
&\qquad gpu\_type\_name, memory\_gb\},\\
\projlabel{Intersection:} &\quad \{request\_id\}^{+} \cap Q = Q.
\end{align*}

\begin{center}
    \textbf{Projected minimal basis:}
\end{center}
\begin{align*}
request\_id &\fd submitted\_at & request\_id &\fd input\_size\\
request\_id &\fd status & request\_id &\fd response\_time\_ms\\
request\_id &\fd request\_region\_id & request\_id &\fd model\_id\\
request\_id &\fd machine\_id.&&
\end{align*}

\code{request\_region\_id} is stored as \code{region\_id} in the table as mentioned above.

\paragraph{CompatibleWith.}
\begin{align*}
\projlabel{Relation:} &\quad CW = \{model\_id, gpu\_type\_name\},\\
\projlabel{Closures:} &\quad \{model\_id\}^{+} \cap CW = \{model\_id\},\\
&\quad \{gpu\_type\_name\}^{+} \cap CW = \{gpu\_type\_name\},\\
\projlabel{Intersection:} &\quad \{model\_id, gpu\_type\_name\}^{+} \cap CW = CW.
\end{align*}
\noindent\textbf{Projected minimal basis:} No non-trivial FDs.

\subsection{Projected Minimal Bases Summary}

\begin{center}
\begin{tabular}{>{\raggedright\arraybackslash}p{0.24\textwidth} >{\raggedright\arraybackslash}p{0.67\textwidth}}
\toprule
\textbf{Relation} & \textbf{Projected minimal basis}\\
\midrule
Region & $region\_id \fd region\_name$, $region\_id \fd country$\\
GPUType & No non-trivial FDs\\
SpecificGPUType & $gpu\_type\_name \fd memory\_gb$\\
AIModel & $model\_id \fd model\_name$\\
GPUServer & $machine\_id \fd$ each non-key server attribute\\
InferenceRequest & $request\_id \fd$ each non-key/FK request attribute\\
MachineMetric & $\{machine\_id,metric\_id\} \fd$ each metric attribute\\
Hosts & No non-trivial FDs\\
CompatibleWith & No non-trivial FDs\\
\bottomrule
\end{tabular}
\end{center}

\section{Dependency Preservation}

Let $P$ be the union of the projected minimal bases above. After mapping the temporary
region role names back to the corresponding local \code{region\_id} columns, $P$ contains
the same 21 non-trivial dependencies as $F_{min}$:

\[
P = F_{min}
\qquad\Longrightarrow\qquad
P^{+}=F_{min}^{+}.
\]

Every dependency can therefore be checked within a single final relation. No dependency
from $F_{min}$ requires joining relations in order to enforce it, so the current
decomposition is \textbf{dependency preserving} for the identified functional dependencies.

\enlargethispage{3\baselineskip}
\section{Conclusion}

The minimal basis was projected onto every relation using the lecture procedure of
calculating closures, intersecting with the target relation, and simplifying the result.
The union of the projected bases is equivalent to $F_{min}$, so the updated Sprint Two
schema preserves the identified functional dependencies.

\end{document}
